\documentclass[11pt,a4paper,]{article}
\usepackage{lmodern}

\usepackage{amssymb,amsmath}
\usepackage{ifxetex,ifluatex}
\usepackage{fixltx2e} % provides \textsubscript
\ifnum 0\ifxetex 1\fi\ifluatex 1\fi=0 % if pdftex
  \usepackage[T1]{fontenc}
  \usepackage[utf8]{inputenc}
\else % if luatex or xelatex
  \usepackage{unicode-math}
  \defaultfontfeatures{Ligatures=TeX,Scale=MatchLowercase}
\fi
% use upquote if available, for straight quotes in verbatim environments
\IfFileExists{upquote.sty}{\usepackage{upquote}}{}
% use microtype if available
\IfFileExists{microtype.sty}{%
\usepackage[]{microtype}
\UseMicrotypeSet[protrusion]{basicmath} % disable protrusion for tt fonts
}{}
\PassOptionsToPackage{hyphens}{url} % url is loaded by hyperref
\usepackage[unicode=true]{hyperref}
\PassOptionsToPackage{usenames,dvipsnames}{color} % color is loaded by hyperref
\hypersetup{
            pdftitle={Forecasting interrupted time series},
            pdfkeywords={blah, blah},
            colorlinks=true,
            linkcolor=blue,
            citecolor=Blue,
            urlcolor=Blue,
            breaklinks=true}
\urlstyle{same}  % don't use monospace font for urls
\usepackage{geometry}
\geometry{left=2.5cm,right=2.5cm,top=2.5cm,bottom=2.5cm}
\usepackage[style=authoryear-comp,]{biblatex}
\addbibresource{references.bib}
\usepackage{longtable,booktabs}
% Fix footnotes in tables (requires footnote package)
\IfFileExists{footnote.sty}{\usepackage{footnote}\makesavenoteenv{long table}}{}
\usepackage{graphicx,grffile}
\makeatletter
\def\maxwidth{\ifdim\Gin@nat@width>\linewidth\linewidth\else\Gin@nat@width\fi}
\def\maxheight{\ifdim\Gin@nat@height>\textheight\textheight\else\Gin@nat@height\fi}
\makeatother
% Scale images if necessary, so that they will not overflow the page
% margins by default, and it is still possible to overwrite the defaults
% using explicit options in \includegraphics[width, height, ...]{}
\setkeys{Gin}{width=\maxwidth,height=\maxheight,keepaspectratio}
\IfFileExists{parskip.sty}{%
\usepackage{parskip}
}{% else
\setlength{\parindent}{0pt}
\setlength{\parskip}{6pt plus 2pt minus 1pt}
}
\setlength{\emergencystretch}{3em}  % prevent overfull lines
\providecommand{\tightlist}{%
  \setlength{\itemsep}{0pt}\setlength{\parskip}{0pt}}
\setcounter{secnumdepth}{5}

% set default figure placement to htbp
\makeatletter
\def\fps@figure{htbp}
\makeatother


\title{Forecasting interrupted time series}

%% MONASH STUFF

%% CAPTIONS
\RequirePackage{caption}
\DeclareCaptionStyle{italic}[justification=centering]
 {labelfont={bf},textfont={it},labelsep=colon}
\captionsetup[figure]{style=italic,format=hang,singlelinecheck=true}
\captionsetup[table]{style=italic,format=hang,singlelinecheck=true}

%% FONT
\RequirePackage{bera}
\RequirePackage[charter,expert]{mathdesign}
\RequirePackage[scale=0.9]{sourcecodepro}
\RequirePackage[lf,t]{FiraSans}

%% HEADERS AND FOOTERS
\RequirePackage{fancyhdr}
\pagestyle{fancy}
\rfoot{\Large\sffamily\raisebox{-0.1cm}{\textbf{\thepage}}}
\makeatletter
\lhead{\textsf{\expandafter{\@title}}}
\makeatother
\rhead{}
\cfoot{}
\setlength{\headheight}{15pt}
\renewcommand{\headrulewidth}{0.4pt}
\renewcommand{\footrulewidth}{0.4pt}
\fancypagestyle{plain}{%
\fancyhf{} % clear all header and footer fields
\fancyfoot[C]{\sffamily\thepage} % except the center
\renewcommand{\headrulewidth}{0pt}
\renewcommand{\footrulewidth}{0pt}}

%% MATHS
\RequirePackage{bm,amsmath}
\allowdisplaybreaks

%% GRAPHICS
\RequirePackage{graphicx}
\setcounter{topnumber}{2}
\setcounter{bottomnumber}{2}
\setcounter{totalnumber}{4}
\renewcommand{\topfraction}{0.85}
\renewcommand{\bottomfraction}{0.85}
\renewcommand{\textfraction}{0.15}
\renewcommand{\floatpagefraction}{0.8}

%\RequirePackage[section]{placeins}

%% SECTION TITLES
\RequirePackage[compact,sf,bf]{titlesec}
\titleformat{\section}[block]
  {\fontsize{15}{17}\bfseries\sffamily}
  {\thesection}
  {0.4em}{}
\titleformat{\subsection}[block]
  {\fontsize{12}{14}\bfseries\sffamily}
  {\thesubsection}
  {0.4em}{}
\titlespacing{\section}{0pt}{*5}{*1}
\titlespacing{\subsection}{0pt}{*2}{*0.2}


%% TITLE PAGE
\def\Date{\number\day}
\def\Month{\ifcase\month\or
 January\or February\or March\or April\or May\or June\or
 July\or August\or September\or October\or November\or December\fi}
\def\Year{\number\year}

\makeatletter
\def\wp#1{\gdef\@wp{#1}}\def\@wp{??/??}
\def\jel#1{\gdef\@jel{#1}}\def\@jel{??}
\def\showjel{{\large\textsf{\textbf{JEL classification:}}~\@jel}}
\def\nojel{\def\showjel{}}
\def\addresses#1{\gdef\@addresses{#1}}\def\@addresses{??}
\def\cover{{\sffamily\setcounter{page}{0}
        \thispagestyle{empty}
        \placefig{2}{1.5}{width=5cm}{_extensions/numbats/wp/monash2}
        \placefig{16.9}{1.5}{width=2.1cm}{_extensions/numbats/wp/MBSportrait}
        \begin{textblock}{4}(16.9,4)ISSN 1440-771X\end{textblock}
        \begin{textblock}{7}(12.7,27.9)\hfill
        \includegraphics[height=0.7cm]{_extensions/numbats/wp/AACSB}~~~
        \includegraphics[height=0.7cm]{_extensions/numbats/wp/EQUIS}~~~
        \includegraphics[height=0.7cm]{_extensions/numbats/wp/AMBA}
        \end{textblock}
        \vspace*{2cm}
        \begin{center}\Large
        Department of Econometrics and Business Statistics\\[.5cm]
        \footnotesize http://monash.edu/business/ebs/research/publications
        \end{center}\vspace{2cm}
        \begin{center}
        \fbox{\parbox{14cm}{\begin{onehalfspace}\centering\Huge\vspace*{0.3cm}
                \textsf{\textbf{\expandafter{\@title}}}\vspace{1cm}\par
                \LARGE\@author\end{onehalfspace}
        }}
        \end{center}
        \vfill
                \begin{center}\Large
                \Month~\Year\\[1cm]
                Working Paper \@wp
        \end{center}\vspace*{2cm}}}
\def\pageone{{\sffamily\setstretch{1}%
        \thispagestyle{empty}%
        \vbox to \textheight{%
        \raggedright\baselineskip=1.2cm
     {\fontsize{24.88}{30}\sffamily\textbf{\expandafter{\@title}}}
        \vspace{2cm}\par
        \hspace{1cm}\parbox{14cm}{\sffamily\large\@addresses}\vspace{1cm}\vfill
        \hspace{1cm}{\large\Date~\Month~\Year}\\[1cm]
        \hspace{1cm}\showjel\vss}}}
\def\blindtitle{{\sffamily
     \thispagestyle{plain}\raggedright\baselineskip=1.2cm
     {\fontsize{24.88}{30}\sffamily\textbf{\expandafter{\@title}}}\vspace{1cm}\par
        }}
\def\titlepage{{\cover\newpage\pageone\newpage\blindtitle}}

\def\blind{\def\titlepage{{\blindtitle}}\let\maketitle\blindtitle}
\def\titlepageonly{\def\titlepage{{\pageone\end{document}}}}
\def\nocover{\def\titlepage{{\pageone\newpage\blindtitle}}\let\maketitle\titlepage}
\let\maketitle\titlepage
\makeatother

%% SPACING
\RequirePackage{setspace}
\spacing{1.5}

%% LINE AND PAGE BREAKING
\sloppy
\clubpenalty = 10000
\widowpenalty = 10000
\brokenpenalty = 10000
\RequirePackage{microtype}

%% PARAGRAPH BREAKS
\setlength{\parskip}{1.4ex}
\setlength{\parindent}{0em}

%% HYPERLINKS
\RequirePackage{xcolor} % Needed for links
\definecolor{darkblue}{rgb}{0,0,.6}
\RequirePackage{url}

\makeatletter
\@ifpackageloaded{hyperref}{}{\RequirePackage{hyperref}}
\makeatother
\hypersetup{
     citecolor=0 0 0,
     breaklinks=true,
     bookmarksopen=true,
     bookmarksnumbered=true,
     linkcolor=darkblue,
     urlcolor=blue,
     citecolor=darkblue,
     colorlinks=true}

%% KEYWORDS
\newenvironment{keywords}{\par\vspace{0.5cm}\noindent{\sffamily\textbf{Keywords:}}}{\vspace{0.25cm}\par\hrule\vspace{0.5cm}\par}

%% ABSTRACT
\renewenvironment{abstract}{\begin{minipage}{\textwidth}\parskip=1.4ex\noindent
\hrule\vspace{0.1cm}\par{\sffamily\textbf{\abstractname}}\newline}
  {\end{minipage}}


\usepackage[T1]{fontenc}
\usepackage[utf8]{inputenc}

\usepackage[showonlyrefs]{mathtools}
\usepackage[no-weekday]{eukdate}

%% BIBLIOGRAPHY

\makeatletter
\@ifpackageloaded{biblatex}{}{\usepackage[style=authoryear-comp, backend=biber, natbib=true]{biblatex}}
\makeatother
\ExecuteBibliographyOptions{bibencoding=utf8,minnames=1,maxnames=3, maxbibnames=99,dashed=false,terseinits=true,giveninits=true,uniquename=false,uniquelist=false,doi=false, isbn=false,url=true,sortcites=false, date=year}

\DeclareFieldFormat{url}{\texttt{\url{#1}}}
\DeclareFieldFormat[article]{pages}{#1}
\DeclareFieldFormat[inproceedings]{pages}{\lowercase{pp.}#1}
\DeclareFieldFormat[incollection]{pages}{\lowercase{pp.}#1}
\DeclareFieldFormat[article]{volume}{\mkbibbold{#1}}
\DeclareFieldFormat[article]{number}{\mkbibparens{#1}}
\DeclareFieldFormat[article]{title}{\MakeCapital{#1}}
\DeclareFieldFormat[inproceedings]{title}{#1}
\DeclareFieldFormat{shorthandwidth}{#1}
% No dot before number of articles
\usepackage{xpatch}
\xpatchbibmacro{volume+number+eid}{\setunit*{\adddot}}{}{}{}
% Remove In: for an article.
\renewbibmacro{in:}{%
  \ifentrytype{article}{}{%
  \printtext{\bibstring{in}\intitlepunct}}}

\makeatletter
\DeclareDelimFormat[cbx@textcite]{nameyeardelim}{\addspace}
\makeatother
\renewcommand*{\finalnamedelim}{%
  %\ifnumgreater{\value{liststop}}{2}{\finalandcomma}{}% there really should be no funny Oxford comma business here
  \addspace\&\space}

\wp{no/yr}
\jel{C10,C14,C22}

\RequirePackage[absolute,overlay]{textpos}
\setlength{\TPHorizModule}{1cm}
\setlength{\TPVertModule}{1cm}
\def\placefig#1#2#3#4{\begin{textblock}{.1}(#1,#2)\rlap{\includegraphics[#3]{#4}}\end{textblock}}


\nocover

\author{Bahman~Rostami-Tabar, Rob J~Hyndman}
\addresses{\textbf{Bahman Rostami-Tabar}\newline
Cardiff Business School\newline
Wales CF10 3EU\newline
United Kingdom\newline
{Email: rostami-tabarb@cardiff.ac.uk}\newline Corresponding author\newline\\[0.5cm]
\textbf{Rob J Hyndman}\newline
Department of Econometrics \& Business Statistics\newline
Monash University\newline
Clayton VIC 3800\newline
Australia\newline
{Email: Rob.Hyndman@monash.edu}\\[0.5cm]
}

\date{\sf\Date~\Month~\Year}
\makeatletter
 \lfoot{\sf Rostami-Tabar, Hyndman: \@date}
\makeatother

\usepackage[section]{placeins}
\makeatletter
\@ifpackageloaded{caption}{}{\usepackage{caption}}
\AtBeginDocument{%
\ifdefined\contentsname
  \renewcommand*\contentsname{Table of contents}
\else
  \newcommand\contentsname{Table of contents}
\fi
\ifdefined\listfigurename
  \renewcommand*\listfigurename{List of Figures}
\else
  \newcommand\listfigurename{List of Figures}
\fi
\ifdefined\listtablename
  \renewcommand*\listtablename{List of Tables}
\else
  \newcommand\listtablename{List of Tables}
\fi
\ifdefined\figurename
  \renewcommand*\figurename{Figure}
\else
  \newcommand\figurename{Figure}
\fi
\ifdefined\tablename
  \renewcommand*\tablename{Table}
\else
  \newcommand\tablename{Table}
\fi
}
\@ifpackageloaded{float}{}{\usepackage{float}}
\floatstyle{ruled}
\@ifundefined{c@chapter}{\newfloat{codelisting}{h}{lop}}{\newfloat{codelisting}{h}{lop}[chapter]}
\floatname{codelisting}{Listing}
\newcommand*\listoflistings{\listof{codelisting}{List of Listings}}
\makeatother
\makeatletter
\makeatother
\makeatletter
\@ifpackageloaded{caption}{}{\usepackage{caption}}
\@ifpackageloaded{subcaption}{}{\usepackage{subcaption}}
\makeatother

% Adjust headwidth in case user has changed geometry in header-includes
\renewcommand{\headwidth}{\textwidth}

\begin{document}
\maketitle
\begin{abstract}
A brief summary of our ideas
\end{abstract}
\begin{keywords}
blah, blah
\end{keywords}

\section{Introduction}\label{introduction}

Time series are sometimes interrupted by unusual events; for example, a
natural disaster may occur, or a there may be a temporary policy change.
Many forecasters have faced this issue recently with the COVID-19
pandemic, where historical patterns have been severely disrupted due to
lockdowns and other restrictions. In this paper, we consider the problem
of forecasting during such an event, and after it has occurred.

This is a different problem from change point detection. In the
situations we consider, we know that a change has occurred, and we want
to forecast the future after the change. Change point detection is about
identifying \emph{when} a change has occurred.

The changes in the series as a result of the disruption may be
relatively simple; for example a left shift at the start of the
disruption, and another at the end of the disruption. Or they may be
more complex, with changes to the seasonal patterns, and changes to the
level, which evolve over time. In this paper, we consider a range of
models that can be used to handle such changes, and compare their
performance on some real data sets.

Time series models usually assume that the data evolve in the future in
a similar way to how they have evolved in the past. But a big event
(such as COVID-19) can result in a future that is different from the
past, at least in the short term.

\section{Handling interruptions when
forecasting}\label{handling-interruptions-when-forecasting}

We consider several possible ways to handle interruptions when
forecasting.

\subsection{Use a highly adaptive
model}\label{use-a-highly-adaptive-model}

Highly adaptive models can adjust to the interruption as it is
happening, and will therefore be able to approximate the data generating
process relatively well. For example, an ETS model with large smoothing
parameters will be able to adjust to the interruption relatively
quickly. This has the advantage of being a very simple solution that is
easy to implement and fast to compute.

However, the prediction intervals will be large, because the model will
have heavily discounted past data. In fact, the model will largely
forget the past data other than the most recent observations, so there
is no memory of the seasonal patterns and other dynamics that were
present before the interruption. Consequently, the approach works best
if there is no assumption that the post-interruption period will be
similar to the pre-interruption period.

\subsection{Use a dynamic regression model with intervention
covariates}\label{use-a-dynamic-regression-model-with-intervention-covariates}

A dynamic regression model with intervention covariates can be used to
model the interruption explicitly. For example, if the intervention
involves a simple level shift, with a reverse level shift at the end of
the intervention, we can use a dummy variable to indicate the
interruption period, and allow the model to adjust to the interruption.
More complicated interventions can be handled by using more covariates.

This has the advantage of retaining the memory of the past, and so the
seasonal patterns and other dynamics will be retained. However, the
model will assume that the post-interruption period will be similar to
the pre-interruption period, and so the prediction intervals may be be
too narrow.

Advantages: retains full memory of the past, and allows the change
period to be effectively modelled provided you choose the covariates
well.

Disadvantages: requires a lot of thought to choose the covariates well.
Assumes that the post-pandemic period will be similar to the
pre-pandemic period.

\subsection{Treat the covid period as missing and use a model that
handles missing
values}\label{treat-the-covid-period-as-missing-and-use-a-model-that-handles-missing-values}

Many time series models will handle periods of missing values. So the
problematic observations that occur during the period of disruption can
be set to missing, and the model should continue to produce forecasts as
if the interruption had not occurred. Of course, the forecasts will not
be accurate for the period of disruption, but they can be interpreted as
``what might have been''. This solution requires a judgement to be made
about when the disruption has begun, and when normality resumes.

Because no information is retained during the disruption, the prediction
intervals will become large during the disruption, and after the
disruption they will remain large until the model has enough data to
estimate the forecast distribution more accurately.

\subsection{Estimate what might have been and adjust the
data}\label{estimate-what-might-have-been-and-adjust-the-data}

A fourth solution is to estimate what might have been during the period
of disruption, and then use the adjusted data to fit a model. The
estimates could be made using any convenient method. One approach would
be to set the estimates to equal the forecasts made using only
pre-interruption data. This then becomes almost equal to the previous
solution except that the prediction intervals will be narrower because
the estimation uncertainty has not been taken into account.

\subsection{Over sampling (weighting)}\label{over-sampling-weighting}

Oversampling is a technique used in machine learning to address class
imbalance within a dataset. In this context, the goal is to ensure that
the model is more influenced by the patterns observed before the
disruption period while still accounting for the observations during the
disruption period.

Using this weighting approach, we can leverage the insights from the
pre-disruptive period while still accounting for the disruptive period's
influence. It allows the model to better capture the underlying patterns
of the time series data, considering both periods, and potentially
producing more accurate forecasts.

Advantages:

Preserves pre-event patterns: By assigning higher weights to the
pre-disruption time series data, the model is encouraged to focus more
on the patterns observed before the disruptive event. This can be
beneficial in capturing the underlying dynamics of the time series that
might be essential for accurate forecasting in the post-disruption era.

Robustness to disruption impact: The weighting approach allows the model
to still take the disruption period into account, even if it's
considered less influential. This ensures that the model is not
completely blind to the disruption's effects and can adapt to changes or
shifts in the time series caused by the event.

Easy to Implement: ??

Interpretability?:

Disadvantages:

\subsection{Ensemble models}\label{ensemble-models}

\section{Examples}\label{examples}

Let's consider some examples

\subsection{Australian tourism}\label{australian-tourism}

The Australian tourism data set is a monthly time series showing the
number of short-term overseas visitors to Australia. The data
\autocite{tourismdata} are available from January 2000 to May 2023, and
are shown in Figure~\ref{fig-tourism-plot1}. As the borders closed in
March 2020, the number of visitors to Australia dropped to near zero,
and remained there until towards the end of 2021. The borders officially
reopened on 21 February 2022, although it seems visitors began to arrive
earlier than that.

\begin{figure}

\centering{

\includegraphics{fits_files/figure-pdf/fig-tourism-plot1-1.pdf}

}

\caption{\label{fig-tourism-plot1}Short-term visitor arrivals to
Australia (monthly): Jan 2000 -- May 2023.}

\end{figure}%

We apply the four solutions discussed in the previous section to these
data, making 12 month forecasts at the end of each year from 2019 to
2022. We fit ETS, ARIMA and dynamic ARIMA models to the data
\autocite{fpp3}, first applying a log transformation to ensure the
resulting forecasts are positive. This is possible, because the
observations never reach exactly zero, with the smallest number of
visitors per month equal to 2250 in April 2020. For all forecasts, we
also show 90\% prediction intervals.

\begin{figure}[t]

\centering{

\includegraphics{fits_files/figure-pdf/fig-tsol1-plot-1.pdf}

}

\caption{\label{fig-tsol1-plot}Solution 1. Forecasts from ETS and ARIMA
models. Neither works particularly well for disruptions of this
magnitude.}

\end{figure}%

In Figure~\ref{fig-tsol1-plot}, we show forecasts by applying ETS and
ARIMA models to the data. ETS, in particular, is well-known to be
relatively adaptive to changes in the series, and this is evident in
these forecasts. The forecasts for 2020 were made using data before
COVID-19 had any affect, and so they show similar patterns to the past.
The forecasts for 2021 were made after 9 months of very low levels of
arrivals, and both models show forecasts consistent with the recent
history. The use of logarithms is particularly important here as the
variance is much smaller during 2020 than previously, but after taking
logarithms, the variance is more stable over time. The forecasts made at
the start of 2022 have struggled to detect the small increase in traffic
at the end of 2021, and both models have forecast relatively flat
trajectories as a result, although the prediction intervals are wide
indicating model uncertainty. Finally, the forecasts made at the end of
2022 have captured the increasing trend, but ETS is much closer to the
reality, adapting more quickly to the changing patterns. Again, the wide
prediction intervals indicate a high level of uncertainty. It is
possible to make ETS more adaptive to changes in the data by increasing
the value of the smoothing parameters. For example, a high value of
\(\beta\) (the smoothing parameter for the slope), will result in
changes of trend being incorporated into the forecasts more quickly, at
the risk of over-reacting to noise in the data, and increasing the size
of the prediction intervals even more.

\begin{figure}[t]

\centering{

\includegraphics{fits_files/figure-pdf/fig-tsol2-plot-1.pdf}

}

\caption{\label{fig-tsol2-plot}Solution 2. We use a left shift from
March 2020 to October 2022, and a ramp from October 2021 to October
2022. After October 2022, there are no non-zero predictors.}

\end{figure}%

Figure~\ref{fig-tsol2-plot} shows forecasts obtained using a dynamic
regression model (i.e., a regression with ARIMA errors) using two
intervention variables: a level shift from March 2022 to November 2022,
and a ramp from October 2021 to November 2022. It is evident that the
level shift variable has worked well, giving relatively good forecasts
for 2021. However, the forecasts for 2022 are particularly poor, because
the ramp slope has been greatly overestimated, as it was based on only
three observations (Oct -- Dec 2021). The forecasts for 2023 are much
better, and the relatively large prediction intervals are appropriate
given the uncertainty in the industry at the end of 2022.

\begin{figure}[t]

\centering{

\includegraphics{fits_files/figure-pdf/fig-tsol3-plot-1.pdf}

}

\caption{\label{fig-tsol3-plot}Solution 3. The period from March 2020 to
October 2022 is set to missing. So the first three years of forecasts
show what might have been. The fourth set of forecasts uses observations
from November and December 2022, and so the forecasts have been adjusted
downwards.}

\end{figure}%

The third solution involved setting the observations during the
disruption period to missing, and then fitting an ARIMA model to the
series. Consequently, the first three years of forecasts show what might
have been without the COVID-19 pandemic, based on the history to the end
of 2019. The final forecasts for 2023 use the data to the end of 2019,
and the two observations in November and December 2022. These are much
better, but the prediction intervals are too narrow, as there hasn't
been sufficient recent data.

\begin{figure}[t]

\centering{

\includegraphics{fits_files/figure-pdf/fig-tsol4-plot-1.pdf}

}

\caption{\label{fig-tsol4-plot}Solution 4. We replace the observations
from March 2020 to October 2022 with the average of the same month in
the three years prior to March 2020.}

\end{figure}%

Finally, we show in Figure~\ref{fig-tsol4-plot} the forecasts obtained
using solution 4. Here we have replaced observations between March 2020
and October 2022 with estimates based on the average of the same month
in the three years prior to March 2020. The resulting forecasts are
similar to those from solution 3, but with narrower prediction intervals
because the model is (falsely) assuming that the ``observations'' during
the disruption are real.

\FloatBarrier

\subsection{Example: Ambulance
attendances}\label{example-ambulance-attendances}

\includegraphics{fits_files/figure-pdf/plot-data-1.pdf}

\includegraphics{fits_files/figure-pdf/plot-data-2.pdf}

\includegraphics{fits_files/figure-pdf/unnamed-chunk-8-1.pdf}

\includegraphics{fits_files/figure-pdf/forecast-1.pdf}

\includegraphics{fits_files/figure-pdf/forecast-2.pdf}

\includegraphics{fits_files/figure-pdf/forecast-3.pdf}

\printbibliography

\end{document}
